\subsection{Flash--Matching Effects on Resolution}\label{sec:FLASH_MATCHING_EFFECT}
As already explained, the flash--matching step in the reconstruction process allows to associate a \ac{TPC} clusters with the corresponding \ac{PDS} flash and hence obtain the correct drift coordinate of the interaction. In the context of energy reconstruction, this has an impact on the drift attenuation correction on a cluster-by-cluster basis. 

\begin{figure}
    \centering
    \includegraphics[width=.49\linewidth]{Images/hd_1x2x6_lateralAPA_Neutrino_SolarEnergy_Resolution.png}
    \includegraphics[width=.49\linewidth]{Images/vd_1x8x14_3view_30deg_nominal_Neutrino_SolarEnergy_Resolution.png}
    \caption{(left) \ac{HD} (right) \ac{VD}. Neutrino energy resolution with and without flash--matching. The dashed lines represent the ideal case where the drift coordinate is taken from truth information.\label{fig:FLASH_MATCHING_EFFECT}}
\end{figure}

As can be seen, in~\cref{fig:FLASH_MATCHING_EFFECT}, the flash--matching step has a residual effect on the energy resolution. The difference between the flash--matched and non-flash--matched cases are small, indicating that the inefficiencies in the energy reconstruction are dominated by other factors, such as clustering inefficiencies and intrinsic energy losses due to nuclear effects and undetected particles. The dashed lines represent the ideal case where the drift coordinate is taken from truth information, showing that even in this idealized scenario, the resolution does not significantly improve.